\section{B-Spline curves}

%% --------------------------------------------------------------------------------------------- %%
%% --------------------------------------------------------------------------------------------- %%
%% --------------------------------------------------------------------------------------------- %%
\subsection{Definition}

A B-Spline curve, shorthand for basis spline curve, is a parametric curve defined by:
\begin{align}
    \label{eq:def_bspline_curve}
    \mathbf{C}(u)=\sum_{i=0}^{n} N_{i,p}(u)\, \mathbf{P}_{i} \quad \quad 0 \leq u \leq 1
\end{align}

where $p$ is the order of the curve, the coefficients $\mathbf{P}_{i}$ are called control points, and $N_{i,p}$ are the basis functions defined on the non-decreasing knot vector $U$:
\begin{align}
    U = [u_0,\, \dots,\, u_{r}] \in \mathds{R}^{r+1} \quad \text{with} \quad r = n + p + 1
\end{align}

The B-Spline basis functions are given by the recursive relation:
\begin{align}
    N_{i,0}(u) =  \begin{cases}
    1& \mathrm{if} \quad u_{i}\leq u<u_{i+1}\\
    0& \mathrm{otherwise}
    \end{cases}
\end{align}
\begin{align}
 N_{i,p}(u)= \frac {u-u_{i}}{u_{i+p}-u_{i}} N_{i,p-1}(u) + \frac {u_{i+p+1}-u}{u_{i+p+1}-u_{i+1}} N_{i+1,p-1}(u)
\end{align}

This recursive relation can yield the $0/0$ quotient which is defined to be zero by convention.



%% --------------------------------------------------------------------------------------------- %%
%% --------------------------------------------------------------------------------------------- %%
%% --------------------------------------------------------------------------------------------- %%
\subsection{Mathematical properties of B-spline basis functions}
Here is a list of some important properties of the B-Spline basis functions:

\begin{enumerate}

    \item The relation $r = n + p + 1$ holds, where $n+1$ is the number of basis functions and $r+1$ is the number of elements of the knot vector $U$.
    
    \item $N_{i,p}(u)$ is, at most, a polynomial of degree $p$.

    \item Local support:
    \begin{enumerate}
        \item $N_{i,p}(u)=0$ if $u$ is outside the interval $[u_{i}, u_{i+p+1})$. This property is illustrated by the De Boor triangular computation.
        \item In any given knot interval, $[u_{i_{0}}, u_{i_{0}+1})$, at most $(p+1)$ basis functions are nonzero, namely the $N_{i,p}(u)$ with $i_{0}-p \leq i \leq i_{0}$.
    \end{enumerate}

    \item Non-negativity:
        \begin{equation*}
            N_{i,p}(u) \geq 0 \quad \text{for all $i$, $p$, and $u \in [0, 1]$}
        \end{equation*}
        
    \item Partition of unity:
        \begin{equation*}
            \sum_{i=0}^n N_{i,n}(u) = 1 \text{ for all } u \in [0, 1]
        \end{equation*}
    
    In addition, for an arbitrary knot span, $[u_{i_{0}}, u_{i_{0}+1})$ we have that:
        \begin{equation*}
            \sum_{i=i_{0}-p}^{i_{0}} N_{i,n}(u) = 1 \text{ for all } u \in [u_{i_{0}}, u_{i_{0}+1})
        \end{equation*}
    This means that the sum of the non-zero basis-functions of any knot span is unity.
        
    \item Continuity and differentiability:
        \begin{enumerate}
            \item The basis functions are infinitely differentiable in the interior of the knot intervals.
            \item The basis functions are $p-k$ continuously differentiable at a knot with multiplicity $k$.
        \end{enumerate}
        
    \item Extrema:
    
    $N_{i,p}(u)$ attains exactly one maximum value in the interval $u \in [0, 1]$.
    
    \item The first derivative of the basis functions is given by:
        \begin{equation*}
            N'_{i,p}(u) = \frac{\mathrm{d}N_{i,p}}{\mathrm{d}u} =
            p \left( \frac{N_{i,p-1}(u)}{u_{i+p}-u_{i}} - \frac{N_{i+1,p-1}(u)}{u_{i+p+1}-u_{i+1}} \right)
        \end{equation*}
        This is proven by induction after a good deal of algebra.
        This equation can yield the $0/0$ quotient which is defined to be zero by convention.
        
    \item The $k$-th order derivative of the basis functions is given by:
        \begin{equation*}
            N^{(k)}_{i,p}(u) = \frac{\mathrm{d}^{(k)} N_{i,p}}{\mathrm{d}u^{(k)}} =
            p \left( \frac{N^{(k-1)}_{i,p-1}(u)}{u_{i+p}-u_{i}} - \frac{N^{(k-1)}_{i+1,p-1}(u)}{u_{i+p+1}-u_{i+1}} \right)
        \end{equation*}
        This is derived by repeated differentiation.
        This equation can yield the $0/0$ quotient which is defined to be zero by convention.

    \item When the first and last knots have multiplicity $p+1$ the knot vector is given by:
        \begin{align*}
            & U = [\underbrace{u_0,\, u_1,\, \dots,\, u_{p}}_{p+1},\, \underbrace{u_{p+1},\, \dots,\, u_{n}}_{n-p},\, \underbrace{u_{n+1},\, \dots,\, u_{n+p},\, u_{n+p+1}}_{p+1}]
        \end{align*}
        where:
        \begin{align*}
            & u_0 = \dots = u_p = 0 \\
            & u_{n+1} = \dots = u_{n+p+1} = 1
        \end{align*}
    and is called a \textit{clamped knot vector}. Basis functions of clamped knot vectors satisfy two additional properties:
    \begin{enumerate}
            \item $N_{0,p}(u=0) =1$ and $N_{i,p}(u=0)=0$ for $i\neq0$
            \item $N_{n,p}(u=1) =1$ and $N_{i,p}(u=1)=0$ for $i\neq n$
    \end{enumerate}
    In the remainder of this note, all knot vectors are understood to be clamped.
    
    \item If the knot vector is clamped and $p=n$, then the B-Spline basis function reduces to Bernstein polynomials, that is, $N_{i,p}(u)=B_{i,n}(u)$. This is true because the recursive definition of the B-Spline basis functions is reduced to the recursive definition of the Bernstein polynomials when the knot vector given by
    \begin{equation*}
            U = [\underbrace{0,\, \dots,\, 0}_{p+1},\, \underbrace{1,\, \dots,\, 1}_{p+1}]
    \end{equation*}
    
\end{enumerate}



%% --------------------------------------------------------------------------------------------- %%
%% --------------------------------------------------------------------------------------------- %%
%% --------------------------------------------------------------------------------------------- %%
\subsection{Mathematical properties of B-Spline curves}
Here is a list of some important properties of B-Spline curves:

\begin{enumerate}

    \item $\mathbf{C}(u)$ is a piecewise curve and its components are polynomials of, at most, degree $p$.
    
    \item The order $p$, number of control points $n+1$ and number of knots $r+1$ are related according to $r = n + p + 1$.
    
    \item If $n = p$ and the knot vector is clamped then $\mathbf{C}(u)$ is a \Bezier curve.
    
    
    \item Affine invariance:
    
    B-Spline curves are invariant under affine transformations such as rotations, displacements and scalings. This means that one can apply an affine transformation to the B-Spline curve by applying it to its set of control points.
    
    \item Convex hull property:
    
    All the points of a B-Spline curve are contained in the \textit{convex hull} of its control points.
    The convex hull of a set of points $P = \{\mathbf{P}_{0},\, \dots,\, \mathbf{P}_{N}\}$ is denoted by $\mathcal{CH}(P)$ and it is the set of all the convex combinations of points:
    \begin{equation*}
     \mathcal{CH}(P) = \Big\{ \mathbf{C} =  \sum_{k=0}^{N} a_{k} \, \mathbf{P}_{k} \text{ such that } \sum_{k=0}^{N} a_{k} = 1 \text{ and } a_{k}\geq 0 \text{ for } k = 0,\, 1,\, \dots,\, N\Big\}
    \end{equation*}
    The convex hull property follows from the non-negativity and the partition-of-unity properties of the basis functions.
    % The convex hull of a set of points is the minimal convex polytope that contains all the points
    
    \item Strong convex hull property:
    
    If $u \in [u_{i_{0}},\, u_{i_{0}+1})$, then $\mathbf{C}(u)$ is contained in the convex hull of the control points $\boldsymbol{P}_{i}$ with $i_{0}-p \leq i \leq i_{0}$.

    \item Local modification scheme:
    
    Modifying the control point $\mathbf{P}_{i}$ affects $\mathbf{C}(u)$ only in the interval $[u_{i}, u_{i+p+1})$. This property follows from the fact that $N_{i,p}(u)=0$ if $u$ is outside the interval $[u_{i}, u_{i+p+1})$ and it implies that the shape of a B-Spline can be modified locally without changing its shape globally (nice motivation in my handwritten notes)
    
    \item The polygon formed by the set of control points is known as \textit{control polygon}. The control polygon represents a piecewise linear approximation to the B-Spline curve.

    \item Variation diminishing property:
    
    No straight line (or plane in three dimensions) intersects the B-Spline curve more times than it intersects its control polygon. An intuitive explanation of this property is that the B-Spline curve does not wiggle more than its control polygon.
    
    \item Continuity and differentiability:
    \begin{enumerate}
        \item B-Spline curves are infinitely differentiable in the interior of the knot intervals.
        \item B-Spline curves are \textit{at least} $p-k$ continuously differentiable at a knot with multiplicity $k$.
    \end{enumerate}
    
    \item The first and higher order derivatives of a B-Spline curve are given by:
        \begin{align*}
            \frac{\mathrm{d}\mathbf{C}}{\mathrm{d}u} &= \sum_{i=0}^{n}N'_{i,p}(u)\, \mathbf{P}_{i} \\
            \frac{\mathrm{d}^{k}\mathbf{C}}{\mathrm{d}u^{k}} &= \sum_{i=0}^{n}N^{(k)}_{i,p}(u)\, \mathbf{P}_{i}
        \end{align*}
        
        The first derivative of a B-Spline curve is also given by:
        \begin{align*}
            \frac{\mathrm{d}\mathbf{C}}{\mathrm{d}u} &= \sum_{i=0}^{n}N'_{i,p}(u)\, \mathbf{P}_{i}
        \end{align*}
        
        evaluated at a new knot vector $U'$ given by:
        \begin{align*}
            & U' = [\underbrace{0,\, \dots,\, 0}_{p},\, \underbrace{u_{p+1},\, \dots,\, u_{n}}_{n-p},\, \underbrace{1,\, \dots,\, 1}_{p}]
        \end{align*}
        
        Note that the derivative of a B-Spline curve of degree $p$ is a new B-Spline curve of degree $p-1$ with a new set of control points and knot vector.
        
        
        The $k$-th derivative of a B-Spline curve is also given by:
        \begin{align*}
            \frac{\mathrm{d}^k\mathbf{C}}{\mathrm{d}u^k} = \sum_{i=0}^{n-k} N_{i,p-k}(u) \, \mathbf{P}_{i}^{(k)}
        \end{align*}
        where:
        \begin{equation*}
            \mathbf{P}_{i}^{(k)} =  \begin{cases}
                \mathbf{P}_{i} & \text{if $k=0$}\\
                \frac{p-k+1}{u_{i+p+1}-u_{i+k}}\, (\mathbf{P}_{i+1}^{(k-1)} - \mathbf{P}_{i}^{(k)}) &\text{if $k>0$}
                \end{cases}
        \end{equation*}
        
        evaluated at the knot vector $U^{(k)}$ given by:
        \begin{align*}
            & U^{(k)} = [\underbrace{0,\, \dots,\, 0}_{p-k+1},\, \underbrace{u_{p+1},\, \dots,\, u_{n}}_{n-p},\, \underbrace{1,\, \dots,\, 1}_{p-k+1}]
        \end{align*}
    

    \item End point interpolation:
    
    The start and end points of a clampled B-Spline curve coincide with the first and last control points, respectively.
    \begin{align*}
        \mathbf{C}(u=0) &= \mathbf{P}_{0} \\
        \mathbf{C}(u=1) &= \mathbf{P}_{n}
    \end{align*}
    
    
    \item End point tangency:
    
    A clamped B-Spline curve is tangent to the control polygon at the endpoints.
    \begin{align*}
        \frac{\mathrm{d}\mathbf{C}}{\mathrm{d}u}\Bigr|_{u=0} &= \left( \frac{p}{u_{p+1}} \right) (\mathbf{P}_{1} - \mathbf{P}_{0}) \\
        \frac{\mathrm{d}\mathbf{C}}{\mathrm{d}u}\Bigr|_{u=1} &= \left( \frac{p}{1-u_{n}} \right) (\mathbf{P}_{n} - \mathbf{P}_{n-1})
    \end{align*}
    
    
    \item End point curvature:
    
    The second derivative of B-Spline curve at its endpoints in given by: \\
    \resizebox{\linewidth}{!}{
    \begin{minipage}{0.99\linewidth}
    \begin{align*}
        \frac{\mathrm{d}^2\mathbf{C}}{\mathrm{d}u^2}\Bigr|_{u=0} &= \frac{p\, (p-1)}{u_{p+1}} \, \left[ \left(\frac{1}{u_{p+2}} \right) \left(\mathbf{P}_{2} - \mathbf{P}_{0} \right) -\left( \frac{1}{u_{p+1}}  + \frac{1}{u_{p+2}} \right) \left( \mathbf{P}_{1} - \mathbf{P}_{0} \right) \right] \\
        \frac{\mathrm{d}^2\mathbf{C}}{\mathrm{d}u^2}\Bigr|_{u=1} &= \frac{p\, (p-1)}{1-u_{n}} \, \left[ \left(\frac{1}{1-u_{n-1}} \right) \left( \mathbf{P}_{n-2} - \mathbf{P}_{n} \right) -\left( \frac{1}{1-u_{n}}  + \frac{1}{1-u_{n-1}} \right) \left( \mathbf{P}_{n-1} - \mathbf{P}_{n} \right) \right]
    \end{align*}
    \vspace{3pt}
    \end{minipage}
    }

    The curvature of a clamped B-Spline curve at its endpoints is given by:
    \begin{align*}
        \kappa(u=0) &=  \left( \frac{p-1}{p}  \right)  \left( \frac{u_{p+1}}{u_{p+2}}  \right) 
        \frac{\norm{ \left(\mathbf{P}_{2}-\mathbf{P}_{0}\right) \times \left(\mathbf{P}_{1}-\mathbf{P}_{0}\right)} }{\norm{\mathbf{P}_{1}-\mathbf{P}_{0}}^3} \\
        \kappa(u=1) &=  \left( \frac{p-1}{p}  \right)  \left( \frac{1-u_{n}}{1-u_{n-1}}  \right)
        \frac{\norm{ \left(\mathbf{P}_{n-2}-\mathbf{P}_{n}\right) \times \left(\mathbf{P}_{n-1}-\mathbf{P}_{n}\right)} }{\norm{\mathbf{P}_{n-1}-\mathbf{P}_{n}}^3}
    \end{align*}

    
    


    
\end{enumerate}



%% --------------------------------------------------------------------------------------------- %%
%% --------------------------------------------------------------------------------------------- %%
%% --------------------------------------------------------------------------------------------- %%




